\chapter{中断、定时器与内核入口}

\section{原理}

轮询模型的问题是 CPU 必须不断主动检查外设。中断把方向反过来：外设在事件发生时打断 CPU，CPU 保存当前执行现场，跳到中断向量表中登记的处理函数，处理最少必要工作，然后恢复被打断的代码。中断是操作系统从“一个主循环”走向“多个异步事件”的关键。

中断处理函数必须短。原因很具体：中断期间可能屏蔽同级或低级中断；处理太久会增加其他事件延迟；处理函数如果访问复杂共享状态，会和主循环或其他中断产生竞态。因此常见设计是 top half 只读取硬件状态、清中断、把事件放入队列；bottom half 或普通任务稍后完成复杂处理。

\begin{lstlisting}
interrupt uart_rx_isr() {
  byte b = UART_DATA;
  ring_push(&rx_queue, b);
  UART_CLEAR_INTERRUPT();
}

void uart_task_step() {
  while (ring_pop(&rx_queue, &b)) {
    parse_byte(b);
  }
}
\end{lstlisting}

中断引入了共享状态竞争。主循环和中断处理函数可能同时访问 \code{rx\_queue}。在单核单片机上，最简单的临界区是短暂关中断，修改队列指针，再开中断。这个方法有效，但必须短；如果关中断期间做慢操作，就会丢事件或增加控制延迟。

\begin{lstlisting}
disable_interrupts();
bool ok = ring_push(&queue, value);
enable_interrupts();
\end{lstlisting}

这个临界区就是锁的祖先。以后在多核机器上，关本核中断不能阻止另一个核心同时访问同一数据结构，于是需要自旋锁、互斥锁和原子操作。但基本问题没有变：某段代码需要在状态不变量被临时打破时不被并发打断。

\subsection{中断控制器为什么不是可选部件}

若只有一个外设和一个 CPU，外设可以直接把中断线连到 CPU。真实系统有多个外设、多个中断源、多个优先级，甚至多个 CPU。中断控制器负责把混乱的硬件事件整理成 CPU 能处理的 vector：记录 pending 位，应用 mask，比较优先级，选择目标 CPU，并在处理完成后接收 EOI。

\begin{lstlisting}
interrupt_controller:
  pending[source] = device_irq_line
  enabled[source] = software_mask
  priority[source] = configured_priority
  target[source] = cpu_id

deliver():
  candidates = pending & enabled
  source = highest_priority(candidates)
  cpu = target[source]
  cpu.raise_interrupt(vector[source])
\end{lstlisting}

没有中断控制器，OS 就不能可靠地屏蔽某类中断，不能区分中断来源，不能把高吞吐设备分散到多个 CPU，也不能在多核系统中发送 IPI。中断控制器是“外设事件”和“CPU 异常入口”之间的路由层。

\subsection{x86-64 中断控制器：local APIC、IOAPIC、MSI/MSI-X}

x86-64 主线里，中断控制器不是一个单独芯片名，而是一组协作机制。local APIC 属于每个 CPU，用于本地 timer、IPI 和接收 vector；IOAPIC 把传统外部 IRQ 重定向到目标 CPU 和 vector；MSI/MSI-X 让 PCIe 设备通过内存写消息投递中断；x2APIC 扩展了 APIC 编址和访问方式，适合大规模多核系统。

\begin{longtable}{p{0.18\linewidth}p{0.34\linewidth}p{0.36\linewidth}}
\toprule
机制 & 主要用途 & OS 关心点 \\
\midrule
local APIC & 每 CPU 中断、本地 timer、IPI & 多核调度、TLB shootdown、timer tick \\
IOAPIC & 把外部 IRQ 重定向到 CPU/vector & IRQ 路由、触发方式、亲和性 \\
MSI/MSI-X & PCIe 设备通过内存写投递中断 & 多队列设备、vector 分配、隔离 \\
x2APIC & 扩展 APIC 编址和访问方式 & 大规模 CPU、虚拟化、IPI 成本 \\
\bottomrule
\end{longtable}

MSI/MSI-X 值得单独强调。传统中断像一根线；MSI 是设备向一个特殊地址写入消息来触发中断。这样 PCIe 设备可以为多个队列分配多个 vector，让网卡 RX queue 0 中断送到 CPU0，RX queue 1 中断送到 CPU1。多队列网络和 NVMe 的伸缩性很大程度依赖这种硬件能力。

\subsection{定时器和时间片}

定时器中断是操作系统调度的起点。没有定时器，只有任务主动让出 CPU 时，系统才能切换任务；一旦任务死循环，整个系统就失控。定时器按固定频率打断 CPU，内核在中断中增加 tick、检查超时、唤醒等待任务，并在必要时设置“需要调度”的标志。

\begin{lstlisting}
interrupt timer_isr() {
  ticks++;
  for each timer in timer_wheel[current_slot]:
    if (timer.expired(ticks)) {
      make_task_runnable(timer.task);
    }
  if (current_task.time_slice_used++ >= current_task.time_slice) {
    need_reschedule = true;
  }
}
\end{lstlisting}

直接扫描所有定时器的复杂度是 \code{O(n)}。如果系统只有几个任务，这很简单；如果有几十万个超时请求，每个 tick 都扫描全表就不可接受。常见算法包括最小堆、时间轮和分层时间轮。最小堆插入和删除是 \code{O(log n)}，查看最近超时是 \code{O(1)}；时间轮把时间划分成槽，插入通常接近 \code{O(1)}，但精度和最大超时范围受槽设计影响。

\subsection{定时器从硬件计数器到 OS 超时}

硬件定时器最小形态是一个递增或递减计数器，加一个比较寄存器。计数器来自固定频率时钟；当计数器达到比较值时，硬件产生中断。OS 在中断里更新软件时间，并设置下一次比较值。

\begin{lstlisting}
program_next_timer(deadline):
  timer.compare = deadline
  timer.enable_interrupt = true

timer_interrupt():
  now = timer.counter
  run_expired_timers(now)
  next = earliest_deadline()
  program_next_timer(next)
\end{lstlisting}

周期 tick 的做法是每隔固定时间中断一次，例如 1ms 或 10ms。高精度和 tickless 的做法是按最近事件编程下一次中断。前者实现简单，后者减少无意义中断，但要求内核精确维护最近 deadline。无论哪种方式，抢占和超时都依赖硬件在未来把控制权交回内核。

\subsection{异常和系统调用}

中断来自外设，异常来自当前指令执行，例如除零、非法指令、缺页、权限错误。系统调用也是一种受控入口：用户程序主动执行特定指令，请求进入内核。三者共同构成内核入口。

\begin{longtable}{p{0.18\linewidth}p{0.32\linewidth}p{0.38\linewidth}}
\toprule
入口 & 来源 & 典型处理 \\
\midrule
中断 & 外设或定时器 & 读取设备状态，唤醒任务，记录 tick \\
异常 & 当前指令失败 & 发送信号、修复缺页、终止进程 \\
系统调用 & 用户主动请求 & 检查参数和权限，操作内核对象 \\
\bottomrule
\end{longtable}

从这里开始，OS 的边界清楚了：用户代码不能直接操作所有资源，而是通过受控入口进入内核；内核根据当前上下文、权限和对象状态决定是否允许请求。

\subsection{异常入口保存的是诊断证据}

异常不是“出错就跳到内核”这么简单。硬件必须保存足够证据，让内核判断能否修复。以 page fault 为例，内核需要知道 faulting address、访问类型、当前特权级、出错指令地址和错误码。以非法指令为例，内核需要知道是哪条指令，以便发送信号或模拟指令。

\begin{longtable}{p{0.22\linewidth}p{0.34\linewidth}p{0.32\linewidth}}
\toprule
证据 & 用途 & 缺失后果 \\
\midrule
fault RIP & 返回或报告出错指令 & 无法重试或定位错误 \\
fault address & 判断哪个虚拟地址出错 & 无法处理缺页或 SIGSEGV \\
access type & 读、写、执行、用户/内核 & 无法区分 COW、NX、权限错误 \\
RFLAGS/CPL & 中断状态、特权级、条件位 & 返回路径不安全 \\
error code/cause & 架构提供的分类原因 & 只能粗暴终止，无法精细恢复 \\
\bottomrule
\end{longtable}

这也是内核崩溃日志要打印寄存器、RIP、栈和 fault address 的原因。它们不是噪声，而是硬件给出的最低层证据。没有这些证据，调试只能猜。

\subsection{上下文保存}

中断、异常和调度都需要保存上下文。x86-64 的最小上下文包括 RIP、RFLAGS、通用寄存器和 RSP。若要切换任务，还要保存当前任务的内核栈、地址空间标识和调度账本。

\begin{lstlisting}
struct Context {
  uint64_t rip;
  uint64_t rsp;
  uint64_t rflags;
  uint64_t regs[16];
};

void switch_to(Task* prev, Task* next) {
  save_context(&prev->ctx);
  current = next;
  restore_context(&next->ctx);
}
\end{lstlisting}

这段伪代码隐藏了真实汇编细节，但说明了不变量：一个任务被切走时，它的下一条指令位置和栈必须能恢复；一个任务被恢复时，寄存器必须像它从未离开 CPU 一样。上下文保存不正确，任何高级调度策略都没有意义。

\subsection{上下文分两类：异常上下文和调度上下文}

异常上下文由硬件入口和低级入口共同保存，描述“被打断的执行现场”。调度上下文描述“线程被切走后下次如何继续”。二者有关但不完全相同。系统调用和中断进入内核时，需要 trap frame；线程主动睡眠或被调度出去时，需要保存 callee-saved 寄存器、内核栈位置和调度状态。

\begin{longtable}{p{0.2\linewidth}p{0.36\linewidth}p{0.32\linewidth}}
\toprule
上下文 & 保存内容 & 典型时机 \\
\midrule
trap frame & 用户 RIP、RSP、RFLAGS、系统调用参数、错误码 & 中断、异常、系统调用入口 \\
switch context & 内核栈指针、callee-saved 寄存器、返回地址 & scheduler 切换内核线程 \\
地址空间上下文 & 页表根、PCID、TLB 状态 & 切换进程地址空间 \\
浮点/SIMD 上下文 & FPU、向量寄存器、控制字 & 用户使用浮点/SIMD 或信号处理 \\
\bottomrule
\end{longtable}

把这几类上下文混在一起会导致很多误解。例如中断打断用户程序时，内核需要 trap frame 才能返回用户态；调度器在两个都已进入内核的线程之间切换时，通常只需要保存内核执行所需的最小寄存器集合。地址空间切换还要改页表根和处理 TLB。

\subsection{trap frame、switch context 和 stack frame 的区别}

“保存现场”这个说法太粗，会掩盖三种完全不同的结构。trap frame 是硬件入口和汇编入口为异常/中断/系统调用保存的用户或被打断现场；switch context 是调度器在内核线程之间切换时保存的最小恢复信息；stack frame 是普通函数调用在栈上的返回地址、旧帧指针和局部变量布局。

\begin{longtable}{p{0.18\linewidth}p{0.30\linewidth}p{0.40\linewidth}}
\toprule
结构 & 谁创建 & 用来回答什么问题 \\
\midrule
trap frame & CPU 硬件入口 + 低级汇编入口 & 从哪个 RIP/RSP/RFLAGS 返回用户态或被打断点 \\
switch context & 调度器和架构切换代码 & 下次调度到该线程时，从哪个内核位置继续 \\
stack frame & 编译器按 ABI 生成的函数调用代码 & 当前函数返回到哪里，局部变量在哪里 \\
signal frame & 内核在用户栈上构造 & 用户态 signal handler 结束后如何恢复原现场 \\
\bottomrule
\end{longtable}

一个线程在系统调用中睡眠时，通常同时拥有这些结构：用户态曾经有 stack frame；进入内核后有 trap frame；内核 C 函数有自己的 stack frame；调度出去时还要保存 switch context。调度器切换的是“当前内核执行点”，不是直接切换到用户程序入口；返回用户态时才消费 trap frame。

\begin{lstlisting}
user calls read()
  -> user stack frame exists
  -> syscall creates trap frame
  -> kernel functions create kernel stack frames
  -> task sleeps, scheduler saves switch context
  -> later scheduler restores switch context
  -> syscall completes and return path restores trap frame
\end{lstlisting}

这也解释了为什么每个线程需要独立内核栈。线程睡在内核里时，它的内核 stack frames 和 trap frame 都还活着；另一个线程进入内核不能覆盖这段栈。

\section{算法与实现分析}

\subsection{中断入口的最小状态机}

中断入口要做的第一件事不是调用 C 函数，而是保存足够现场，让被打断代码将来能继续。可以把入口写成四段：硬件自动保存、汇编入口补充保存、C 层分发、返回路径恢复。

\begin{lstlisting}
interrupt_entry(vector):
  frame = hardware_saved_frame()
  save_scratch_registers(frame)
  switch_to_kernel_stack_if_needed(frame)
  dispatch_interrupt(vector, frame)
  maybe_preempt_on_return(frame)
  restore_registers(frame)
  iret(frame)
\end{lstlisting}

这个状态机维护两个不变量。第一，任何会被中断处理函数破坏的寄存器，要么已经保存，要么按 ABI 规定不需要保存。第二，返回用户态之前，内核必须处理挂起信号、抢占标志和可能的调度请求。若在中断中直接切换任务，必须保证旧任务的内核栈仍保存完整 trap frame。

\subsection{精确中断和延迟中断}

外部中断通常在指令边界被接收，CPU 完成当前指令的架构效果后进入中断入口。这样内核返回时可以从下一条指令继续。异常则和当前指令绑定：缺页异常通常要返回到同一条指令重试，除零或非法指令可能发送信号。精确性让 OS 能用同一套保存/恢复机制处理多种入口。

\begin{lstlisting}
external interrupt:
  instruction N committed
  interrupt taken before N+1
  saved RIP = address of N+1

page fault:
  instruction N cannot complete
  saved RIP = address of N
  kernel may fix page and retry N
\end{lstlisting}

中断延迟来自多个来源：CPU 临时关中断、硬件优先级屏蔽、长时间持有自旋锁、中断控制器路由、cache miss、下半部积压。实时系统关心最大延迟；通用系统也要避免长时间关中断，因为它会让定时器、网络、块 I/O 和 RCU 都延后。

\subsection{中断控制器与优先级}

真实硬件通常通过中断控制器管理多路中断。控制器至少要完成屏蔽、优先级、确认和结束通知。教学模型可以抽象成 pending bitmap：

\begin{lstlisting}
interrupt_controller_tick():
  pending = pending_bits & enabled_bits
  if pending == 0:
    return
  vector = highest_priority(pending)
  pending_bits.clear(vector)
  cpu_raise_interrupt(vector)

end_of_interrupt(vector):
  controller.eoi(vector)
\end{lstlisting}

若忘记 EOI，控制器可能不再投递后续同类中断；若过早 EOI，同一个设备可能在 handler 尚未处理完共享状态时再次打断。驱动通常先读取设备状态并清设备侧中断源，再向控制器确认完成。顺序错误会导致中断风暴或丢事件。

\subsection{IRQ、vector、handler 和 completion 的区别}

设备 I/O 里还有一组容易混的词。IRQ 是“有中断事件需要处理”的硬件/内核资源；vector 是 CPU 入口编号；handler 是内核注册的处理函数；completion 是某个具体 I/O 请求完成的事实。一个 IRQ 可以对应一个队列，一次 handler 可以处理多个 completion，一个 completion 也可能在轮询模式下被发现而不是靠中断发现。

\begin{longtable}{p{0.18\linewidth}p{0.30\linewidth}p{0.40\linewidth}}
\toprule
概念 & 表示什么 & 常见误解 \\
\midrule
IRQ line/vector & CPU 被通知有事件 & 误以为它等于某个请求完成 \\
handler & 处理该类事件的函数 & 误以为 handler 只能处理一个事件 \\
completion entry & 设备写回的请求结果 & 误以为一定伴随一次中断 \\
wait queue & 等待某个条件的任务集合 & 误以为中断直接唤醒某个用户函数 \\
poll budget & 一次下半部最多处理多少完成 & 误以为只要有中断就会处理完所有包/请求 \\
\bottomrule
\end{longtable}

高吞吐设备通常把中断当成“提醒你去看队列”，而不是“这个请求完成了”的一对一消息。网卡可能一次中断后处理几十个 RX completion；NVMe 可能一个 MSI-X vector 对应一个 completion queue；NAPI 甚至会暂时关闭中断，改用轮询批量处理。这样吞吐更好，但延迟、budget、队列深度和背压都必须进入设计。

\subsection{电平触发、边沿触发和 MSI}

中断触发方式影响 handler 的确认顺序。边沿触发在信号变化瞬间产生事件；若软件没及时处理，可能错过后续边沿。电平触发只要设备中断线保持有效，控制器就认为中断仍 pending；若 handler 没清设备状态，EOI 后会立刻再次进入。MSI/MSI-X 则是设备写一个中断消息，通常不依赖共享物理线。

\begin{longtable}{p{0.18\linewidth}p{0.34\linewidth}p{0.36\linewidth}}
\toprule
触发方式 & 特点 & handler 重点 \\
\midrule
边沿触发 & 变化瞬间产生一次事件 & 不能丢边沿，事件队列要足够 \\
电平触发 & 条件保持则持续 pending & 必须清设备侧原因再 EOI \\
MSI/MSI-X & 设备写消息触发 vector & vector 分配、队列亲和性、屏蔽同步 \\
\bottomrule
\end{longtable}

很多“中断风暴”来自清除顺序错误或设备状态未消费。OS 不能只从 CPU 侧看中断，还要看设备侧 pending 条件是否被真正清掉。

\subsection{top half 与 bottom half 的分工}

中断上半部只做必须立即完成的工作：读取硬件状态、清中断源、把事件入队、唤醒下半部。下半部可以在软中断、工作队列或普通内核线程中运行。

\begin{lstlisting}
irq_top_half(dev):
  status = mmio_read(dev.status)
  mmio_write(dev.ack, status)
  push_event(dev.event_queue, status)
  schedule_bottom_half(dev)

bottom_half(dev):
  while pop_event(dev.event_queue, &event):
    process_event(event)
    wake_waiters(dev.waitq)
\end{lstlisting}

分工的理由是延迟控制。上半部越长，其他中断越晚被处理；下半部越慢，队列越容易积压。因此测试不只看功能，还要看最大关中断时间、队列峰值和事件丢失计数。

\subsection{定时器数据结构选择}

若每个 tick 扫描所有定时器，复杂度是 \code{O(n)}。最小堆适合定时器数量中等、过期时间分散的系统；时间轮适合大量短超时；分层时间轮适合同时覆盖短超时和长超时。

\begin{longtable}{p{0.18\linewidth}p{0.24\linewidth}p{0.24\linewidth}p{0.22\linewidth}}
\toprule
结构 & 插入 & 取最近超时 & 适用场景 \\
\midrule
线性表 & \code{O(1)} & \code{O(n)} & 少量定时器，教学 baseline \\
最小堆 & \code{O(log n)} & \code{O(1)} & 通用超时，精度清楚 \\
单级时间轮 & 接近 \code{O(1)} & 接近 \code{O(1)} & 大量短超时，槽宽可接受 \\
分层时间轮 & 接近 \code{O(1)} & 接近 \code{O(1)} & 长短超时混合 \\
\bottomrule
\end{longtable}

选择定时器结构时要写清精度、最大超时、插入成本和 tick 成本。一个只在 10 个定时器上正确的实现，不能直接推广到 100000 个 socket 超时。

\section{实践}

本章实践建议写一个事件日志。每次进入中断、退出中断、唤醒任务、设置 need\_reschedule，都记录 tick、事件类型和当前任务。然后用日志回答：哪个事件打断了谁，哪个任务被唤醒，什么时候真正运行。

\begin{lstlisting}
tick=100 enter timer_isr current=control
tick=100 wake task=telemetry reason=period
tick=100 set need_reschedule
tick=100 exit timer_isr current=control
tick=100 switch control -> telemetry
\end{lstlisting}

同时实现两个定时器队列版本：全表扫描和最小堆。比较 n 从 10 到 100000 时，每个 tick 的处理成本。这个实验能说明为什么 OS 算法不能只看功能正确，还要看事件规模。

\section{测试}

最低测试清单如下：

\begin{itemize}
\tightlist
\item 中断处理函数只做短路径工作，不调用可能阻塞的函数。
\item 环形队列在主循环和 ISR 同时访问时不破坏 head/tail/size。
\item 定时器 tick 单调增加，超时任务只被唤醒一次。
\item 时间片用尽后设置 reschedule 标志，不能在任意位置破坏当前上下文。
\item 上下文切换后，任务能从被切走的位置继续执行。
\item 异常路径能区分可修复缺页、权限错误和不可恢复错误。
\end{itemize}

这一章的验收结论是：中断和定时器让系统获得外部事件和抢占能力，但也迫使我们认真处理共享状态、上下文保存和入口权限。

\section{源码级补充：x86 IDT、APIC 与 Linux 中断路径}

\subsection{x86 异常入口保存了什么}

x86-64 的中断和异常入口由 IDT 描述。CPU 根据 vector 找到 gate，切换特权级时切到内核栈，硬件压入返回所需状态。某些异常还会压入 error code。汇编入口随后保存通用寄存器，形成内核能理解的寄存器帧。

\begin{lstlisting}
struct TrapFrame {
  // software-saved general registers
  u64 r15, r14, r13, r12, r11, r10, r9, r8;
  u64 rdi, rsi, rbp, rbx, rdx, rcx, rax;

  // hardware-saved frame
  u64 vector;
  u64 error_code;
  u64 rip;
  u64 cs;
  u64 rflags;
  u64 rsp;
  u64 ss;
};
\end{lstlisting}

真实布局会因入口类型和内核版本变化，但读源码时要抓住语义：\code{rip/cs/rflags/rsp/ss} 决定如何返回；error code 帮助 page fault 和保护异常分类；通用寄存器保存系统调用参数和被打断计算状态。

\subsection{IDT、APIC、IOAPIC 和 MSI}

IDT 是 CPU 本地表；APIC/IOAPIC/MSI 解决“设备事件如何送到哪个 CPU 的哪个 vector”。传统外设中断可经 IOAPIC 重定向到本地 APIC；PCIe 设备常用 MSI/MSI-X，通过一次内存写消息触发中断。多核重调度和 TLB shootdown 使用 IPI，由一个 CPU 通过本地 APIC 向另一个 CPU 投递。

\begin{lstlisting}
device event
  -> IOAPIC redirection entry or MSI message
  -> local APIC on target CPU
  -> CPU vector
  -> IDT gate
  -> low-level entry
  -> generic IRQ dispatch
  -> device handler
  -> EOI
\end{lstlisting}

中断亲和性影响性能。网卡 RX 队列中断若全部打到 CPU0，CPU0 会成为瓶颈；把队列中断分散到处理对应网络栈的 CPU，可以减少跨核 cache line 迁移。但过度分散也会让连接状态在 CPU 间跳动。真实系统因此有 irq affinity、RSS/RPS/XPS 等策略。

\subsection{Linux 中断路径阅读点}

读 Linux 中断路径可以从这些位置开始：

\begin{itemize}
\tightlist
\item \filepath{arch/x86/kernel/idt.c}：IDT gate 初始化。
\item \filepath{arch/x86/entry/entry\_64.S}：低级入口、寄存器保存、返回路径。
\item \filepath{kernel/irq/handle.c}：通用 IRQ 分发。
\item \filepath{kernel/irq/manage.c}：\code{request\_irq}、\code{request\_threaded\_irq}。
\item \filepath{kernel/time/timer.c}、\filepath{kernel/time/hrtimer.c}：低精度 timer wheel 与 hrtimer。
\end{itemize}

阅读时不要只看 handler 函数。要同时看谁注册 handler、handler 运行在哪种上下文、是否允许睡眠、谁发送 EOI、错误返回如何统计、设备移除时如何注销并等待 in-flight handler 完成。

\subsection{四种下半部机制}

\begin{longtable}{p{0.2\linewidth}p{0.34\linewidth}p{0.34\linewidth}}
\toprule
机制 & 特点 & 适用场景 \\
\midrule
softirq & 固定类型，常用于网络和块层，高并发 & 极热路径，要求不能随意睡眠 \\
tasklet & 基于 softirq 的延后执行，同一 tasklet 串行 & 旧驱动常见，新代码逐步少用 \\
workqueue & 在内核线程中执行，可睡眠 & 需要分配内存、拿 mutex、访问慢设备 \\
threaded IRQ & 把 IRQ handler 线程化 & 降低硬中断时间，适合较复杂设备处理 \\
\bottomrule
\end{longtable}

选择机制的核心问题是上下文：是否需要睡眠，是否要高吞吐，是否允许同一事件并行，是否要限制 CPU 占用。把会睡眠的代码放进 hard IRQ 或 softirq，是典型内核 bug。

\subsection{hrtimer、timer wheel 与 tickless}

低精度超时适合 timer wheel，例如 TCP 重传、普通超时、周期任务；高精度计时适合 hrtimer，例如纳秒级睡眠、音视频定时、精确 profiler。tickless 则在 CPU 空闲或单任务运行时减少周期 tick，降低功耗和噪声。

\begin{lstlisting}
timer wheel:
  cheap insertion, coarse granularity, many ordinary timeouts

hrtimer:
  ordered by expiry in high-resolution time, higher overhead, precise wakeup

tickless:
  program next hardware timer to nearest event instead of fixed periodic tick
\end{lstlisting}

tickless 不表示没有时间。它表示内核不再无条件固定频率打断 CPU，而是计算下一个必要事件并设置硬件定时器。调度统计、RCU、负载均衡和超时都要适配这种模式。
